\chapter*{Sammendrag}

I norsk akuttmedisin benyttes videostrømming som beslutningsstøtte for \gls{amk}- og \gls{lvs}-operatører. Videoløsningene driftes av \gls{hdo} og baserer seg i dag på to ulike \gls{videomotor}er: AntMedia Server og \gls{nhn} Video. Hver videomotor benytter ulike \gls{api}-er, autentiseringsmekanismer og domenemodeller.

Bacheloroppgaven undersøker i hvilken grad en \gls{provider}-basert \gls{mellomvare}arkitektur kan eksponere ett felles API mot klientene, samtidig som logikk som er spesifikk for hver videomotor isoleres fra kjernen i systemet. Forskningsspørsmålet evalueres langs tre dimensjoner: API-dekning, backend-isolering og ikke-funksjonelle krav. Begrensninger ved heterogene domenemodeller diskuteres som en del av vurderingen.

Prosjektet ble gjennomført som et design- og implementasjonsprosjekt med iterativ arbeidsform. Det resulterende API-et er utviklet i ASP.NET Core med egne provider-implementasjoner for AntMedia Server og NHN Video, \gls{token}-basert autentisering via \gls{keycloak}, containerisert distribusjon med Docker Compose, helsesjekker, strukturert logging og Prometheus-metrikker visualisert i Grafana. Løsningen ble evaluert gjennom kravsporing, automatiserte enhets- og integrasjonstester, et ende-til-ende-testskript mot HDOs testmiljøer, arkitekturvalidering ved tilføyelse av to ekstra providere og en strukturert sluttevaluering fra oppdragsgiver.

Resultatene viser at arkitekturen dekker kjerneoperasjonene oppretting, henting og sletting av videorom på en ensartet måte, og at backend-spesifikk logikk i stor grad lar seg isolere i provider-laget. Samtidig viser arbeidet at full standardisering er vanskelig når videomotorene bygger på ulike domenemodeller, særlig for start og stopp av videostrøm, der semantikken ikke er sammenlignbar mellom de to motorene.

Oppgaven konkluderer med at en provider-basert mellomvare er en egnet tilnærming for HDOs \gls{proof-of-concept}-kontekst. Før eventuell produksjonssetting bør løsningen videreutvikles med strengere autorisasjonsmodell, produksjonsklar tokenvalidering, lasttesting og mer robust håndtering av feiltilstander i eksterne videomotorer.